% Proposed B.90; complete reviewed first-stage argument, awaiting final source review.
\statement{B.90}{Proved}{The first actual relative centralization of the symmetric intersection has surjective fixed-element lifting}
\label{stmt:B90}
Use the based spaces, involution, and actual filling groups of B.89. Put
\[
 \begin{gathered}
 K_s=\ker(\Gamma\to Q)=N_A N_B,\qquad
 M=N_A\cap N_B=[N_A,N_B],\\
 R_A=\ker(G_1\to Q),\qquad R_B=\ker(G_2\to Q).
 \end{gathered}
\]
Then $K_s/M=R_A\times R_B$. Let $\mathscr A,\mathscr B$ be the
sets of connected components of the inverse images of $A,B$ in the
universal cover $\widetilde X$ of the actual filled surface. There is
an isomorphism of abelian groups with involution
\begin{equation}\label{eq:B90-module}
 C:=M/[M,K_s]\cong
 \mathbb Z\bigl[\{(D_A,D_B)\in\mathscr A\times\mathscr B:
                         D_A\cap D_B=\varnothing\}\bigr],
 \qquad F[D_A,D_B]=-[F(D_B),F(D_A)].
\end{equation}
Every index orbit on the right has two elements. In particular
\begin{equation}\label{eq:B90-H1}
 H^1(C_2,C)=0,
 \qquad
 (\Gamma/[M,K_s])^F\twoheadrightarrow
                  (G_1\times_QG_2)^F.
\end{equation}
Let $\mathcal M=\overline M$ and
$\mathcal K_s=\ker(\widehat\Gamma\to\widehat Q)$ be the actual
closed subgroups. With closed commutators in the compact expression,
the actual full completions also satisfy
\begin{equation}\label{eq:B90-compact}
 \bigl(\widehat\Gamma/[\mathcal M,\mathcal K_s]\bigr)^F
       \twoheadrightarrow
       (\widehat G_1\times_{\widehat Q}\widehat G_2)^F.
\end{equation}
The last conclusion is proved by density of the actual discrete fixed
targets. It does not assert a cohomology-vanishing theorem for an
unspecified completion of the coefficient module in
\eqref{eq:B90-module}. Fixed lifting through $M$ itself, and the
projected fixed-source equality, remain unresolved.

\begin{proof}
\emph{B.90.1. The precise inherited source input.}
Let
\[
 S_0=X\cap\{v+w=2\},\qquad S=S_0\setminus A,
 \qquad H_b=\operatorname{im}(\pi_1(S_0)\to Q).
\]
B.87 proves that the displayed source map is injective and
$H_b=Q^F$, with the corresponding equality in full completions.
Write
\[
 I=\operatorname{im}\bigl(
       \ker(\pi_1(S)\to\pi_1(S_0))\longrightarrow R_A\bigr).
\]
The supplied LM4 F.21--F.22 and the actual covering comparisons in
equations B.78.1--3 and B.78.15--17 give the discrete assertion
\begin{equation}\label{eq:B90-normal-generation}
 \langle\!\langle I\rangle\!\rangle_{R_A}=R_A,
 \qquad H_1(I,\mathbb Z)\twoheadrightarrow H_1(R_A,\mathbb Z).
\end{equation}
We spell out why this is an internal, discrete input. LM4 F.21 proves
transitivity of the actual source on each full discrete cross-divisor
coset space, as well as the discrete tubular lift of its divisor
group to the centralizer of the selected meridian. Its factorization
$g=d h c$, with $d$ in the filling kernel, $h$ in the actual source,
and $c$ centralizing that meridian, therefore applies before
completion. Every ambient conjugate of a selected meridian becomes
a conjugate, by the filling kernel itself, of an element of the
source kernel. The selected meridians normally generate the filling
kernel, proving internal normal generation discretely. The actual
epimorphism in LM4 F.22 carries that assertion to the eight-collision
marginal kernel. Both covering characters in B.78 factor through
both ambient and source fillings, so their kernels are literally
unchanged. This gives precisely the $I\subseteq R_A$ just displayed.
Passing to abelianizations proves the second assertion in
\eqref{eq:B90-normal-generation}. No comparison of ambient
metabelian images is substituted for this statement.

\emph{B.90.2. The actual homology calculation.}
Pull $Y,V,Z$ back to $\widetilde X$. The resulting covers
$\widetilde Y,\widetilde V,\widetilde Z$ are connected, because
the actual maps to $Q$ are onto. The one-sided covers have groups
$R_A,R_B$. These groups are free: the two bundle sequences identify
each with the kernel of the actual free-fiber filling $F_{13}\to F_5$.
Since $Y,V$ are aspherical by B.89, the covers
$\widetilde Y,\widetilde V$ are homotopy equivalent to graphs.

The universal cover $\widetilde X$ is contractible. Its punctured
open subset
$\widetilde Y\cup\widetilde V$ is obtained by removing all lifts
of the $64$ transverse nodes. Excision and the relative homology
sequence in real dimension four give
\[
 H_3(\widetilde Y\cup\widetilde V,\mathbb Z)
       =\mathbb Z[Q]^{64},\qquad
 \widetilde H_i(\widetilde Y\cup\widetilde V,\mathbb Z)=0
       \quad(i=1,2).
\]
These are direct sums: singular cycles have compact support and meet
only finitely many of the locally finite lifted nodes. Mayer--Vietoris
therefore gives
\begin{equation}\label{eq:B90-MV}
 H_2(\widetilde Z,\mathbb Z)=\mathbb Z[Q]^{64},
 \qquad
 H_1(\widetilde Z,\mathbb Z)=R_A^{\mathrm{ab}}\oplus R_B^{\mathrm{ab}}.
\end{equation}
The first map takes each local linking three-sphere to the local
linking two-torus, with a consistent orientation. Interchanging the
two open sets changes its Mayer--Vietoris sign.

The Hopf map $H_2(\widetilde Z,\mathbb Z)\to H_2(K_s,\mathbb Z)$
is onto. In the five-term homology sequence for
$1\to M\to K_s\to R_A\times R_B\to1$, the last map on
abelianizations is an isomorphism by \eqref{eq:B90-MV}. Since the
two factors are free, their product has
$H_2=R_A^{\mathrm{ab}}\otimes_{\mathbb Z}R_B^{\mathrm{ab}}$.
Consequently
\begin{equation}\label{eq:B90-tensor-quotient}
 M/[M,K_s]\cong
 \frac{R_A^{\mathrm{ab}}\otimes_{\mathbb Z}R_B^{\mathrm{ab}}}
      {\langle\text{the actual local node tensors}\rangle}.
\end{equation}

Codimension-two excision in the contractible $\widetilde X$ identifies
$R_A^{\mathrm{ab}}$ with the free abelian group on $\mathscr A$,
and similarly for $R_B$. Each lifted component is a lift of an
actual section over the product-pencil base; in particular the
meridian has its positive complex orientation. A local node torus
maps to exactly the basis tensor of its two incident components,
with coefficient one. Repeated intersections of the same two lifted
components merely repeat the same relation. Thus the quotient in
\eqref{eq:B90-tensor-quotient} has the basis of disjoint component
pairs in \eqref{eq:B90-module}. The involution exchanges the two
meridians, so the cross commutator changes sign. This proves the
specified negative-swap action, also agreeing with the
Mayer--Vietoris sign.

\emph{B.90.3. The actual source excludes every singleton orbit.}
Choose the lift $\widetilde S_0$ of $S_0$ through the fixed
basepoint in $\widetilde X$. B.87 makes it the simply connected
cover of $S_0$. The based lift of $F$ fixes it pointwise, by unique
path lifting. Its punctured complement
$\widetilde S=\widetilde S_0\setminus\widetilde A$ is the
covering of $S$ corresponding to the actual source filling kernel;
its image in $R_A$ is $I$.

The group $H_1(\widetilde S,\mathbb Z)$ is generated by the small
meridians of its actual intersection punctures. Each maps to the
single positive meridian basis vector of the lifted $A$ component
which it meets. All these source intersections are transverse.
Equation~\eqref{eq:B90-normal-generation} and the independent
meridian basis therefore force every $D_A\in\mathscr A$ to
meet $\widetilde S_0$. At such a point $x$, one has $F(x)=x$,
so $x\in D_A\cap F(D_A)$.

An index in \eqref{eq:B90-module} can be fixed under the index
permutation only if $D_B=F(D_A)$. The preceding argument excludes
exactly those indices from the disjoint-pair set. Every remaining
orbit has two elements, say $e,e'$ with $F(e)=-e'$ and
$F(e')=-e$. This signed permutation module is free over
$\mathbb Z[C_2]$ after changing the sign of one basis element.
More explicitly, a cocycle on this summand has the form
$a(e+e')=m-F(m)$ with $m=ae$. A cocycle in the direct sum has
finite support, so summing these solutions proves
$H^1(C_2,C)=0$ integrally, without dividing by two.

\emph{B.90.4. Fixed lifting in the discrete quotient.}
Let $D=G_1\times_QG_2$ and $\Gamma_2=\Gamma/[M,K_s]$.
The actual $F$-equivariant sequence is
\[
 1\longrightarrow C\longrightarrow\Gamma_2
                    \longrightarrow D\longrightarrow1.
\]
For $d\in D^F$ and a lift $x$, the right defect $x^{-1}F(x)$
lies in $C$ and is a $C_2$-cocycle. By B.90.3 it equals
$m-F(m)$ in additive notation. Replacing $x$ by $xm$ makes
the right defect zero. This proves the discrete surjection in
\eqref{eq:B90-H1}.

\emph{B.90.5. The actual compact conclusion.}
Identify $G_2$ with $G_1$ by $F$, and let $f:G_1\to Q$ be
the first filling map. Then
\[
 D^F=\operatorname{diag}(f^{-1}(H_b)),\qquad
 \widehat D^{\,F}
   =\operatorname{diag}(\widehat f^{-1}(\overline H_b)).
\]
The second equality uses the actual compact fiber product of B.89
and the full fixed-subgroup equality of B.87.
The discrete $f^{-1}(H_b)$ is dense in the latter compact preimage:
its closure contains $\overline{R_A}=\ker\widehat f$, and its
compact image is all of $\overline H_b$. A closed subgroup
containing that full kernel and projecting onto $\overline H_b$
is the full preimage.

The image of the compact fixed subgroup
$(\widehat{\Gamma_2})^F$ in $\widehat D$ is closed. B.90.4
puts the dense image of $D^F$ inside it, so that image is exactly
$\widehat D^{\,F}$. Finally B.89 identifies
$\overline M=\mathcal M$ with the actual compact intersection,
and density and continuity of commutation give
\[
 \widehat{\Gamma_2}
 =\widehat\Gamma/\overline{[M,K_s]}
 =\widehat\Gamma/[\mathcal M,\mathcal K_s].
\]
This proves \eqref{eq:B90-compact}. It avoids an identification of
the actual completed coefficient kernel with a freely completed
permutation module, and it does not deduce cohomology vanishing
after an arbitrary quotient of such a module.

\emph{B.90.6. Scope and retained proof.}
The result removes the first obstruction after centralizing $M$
relative to $K_s$. Passing to a further nonabelian layer requires
a new argument. In particular neither
$\Gamma^F\to D^F$ nor its compact counterpart is proved onto,
and no projected fixed-source equality is inferred.
The exact source transfer, homology argument, and independent
reviews are retained in
\path{work_A/symmetric_first_fixed_lift_module_16h.md} and the
review files named there. This proof uses the actual source
normal-generation theorem, not a finite approximation to it.
\end{proof}

